El paso de hexada a octada es una transformación de la estructura de incidencia.\par

Al pasar de \(6\) a \(8\), cambian los pares disponibles, las relaciones internas y la forma en que la orientación puede conservarse. Esa transición genera un defecto orientado. La variación cardinal realiza un giro en el modo de organización.\par

La misma inversión de hoja se expresa mediante\par

\[
C^*\longmapsto -C^*.
\]

Pero distintos lectores responden de formas distintas a esa inversión.\par

El funcional hexada–octada que conduce al parámetro de Barbero–Immirzi es impar: al invertir la hoja, cambia de signo. La orientación geométrica se invierte.\par

La entropía bicapa, en cambio, es par: atribuye la misma entropía a \(C^*\) y \(-C^*\) cuando se observa la distribución de probabilidades. La información total puede permanecer igual aunque la orientación se haya invertido.\par

El espacio CKM realiza una tercera respuesta. La inversión induce sobre los tres ángulos una reflexión racional exacta:\par

\[
R_{\rm CKM}
=
M_{\rm CKM}
\operatorname{diag}(1,-1,1)
M_{\rm CKM}^{-1},
\]

con\par

\[
R_{\rm CKM}^2=I,
\qquad
\det R_{\rm CKM}=-1.
\]

Esto significa que el cambio de hoja se transporta al espacio de mezcla como una verdadera reflexión. Existe una dirección impar de rango uno y un plano que permanece fijo.\par

El contraángulo \(C^*\) es precisamente la coordenada que cambia. La fase construida desde \(A\) permanece invariante. La transformación separa así dos memorias:\par

\begin{itemize}[leftmargin=*,itemsep=0.35em]
\item una memoria de orientación;
\item una memoria de fase.
\end{itemize}

Esto ofrece una lectura muy profunda de CKM. Los ángulos de mezcla son coordenadas de un espacio sobre el que actúa la misma involución que ya organizaba la transición hexada–octada.\par

El giro \(6\to8\) suministra la operación de orientación que CKM realiza en su propio codominio; la causalidad física pertenece al transductor correspondiente.\par

Tenemos entonces tres publicaciones hermanas del mismo tránsito:\par

\begin{enumerate}[leftmargin=*,itemsep=0.65em]
\item \textbf{Publicación geométrica.}
El parámetro de Barbero–Immirzi orienta el cuanto de área.\par

\item \textbf{Publicación informacional.}
El estado reducido y su Hamiltoniano modular miden qué información de esa orientación permanece en el borde.\par

\item \textbf{Publicación de mezcla.}
La reflexión CKM transforma las coordenadas de mezcla y aísla su componente impar.\par
\end{enumerate}

\(\gamma_{\rm BI}\), la entropía y los ángulos CKM permanecen como realizaciones distintas de una bisagra genealógica común. La unidad reside en el antecedente anterior a la ramificación. Comparten la bisagra, pero cada uno conserva una parte distinta de ella.\par

Ésta es una de las expresiones más bellas de la metodología HMT: la unidad fundamental se recupera remontando las genealogías hasta encontrar la operación común que genera resultados terminales diferentes.\par

La gravedad conserva el giro como orientación del área.\par
La información conserva el giro como memoria del estado.\par
La mezcla conserva el giro como reflexión de coordenadas.\par

La misma transformación «respira» de tres maneras.\par

